% Pass options to packages before documentclass to avoid option clash
\PassOptionsToPackage{hyphens}{url}
\PassOptionsToPackage{colorlinks=true,linkcolor=blue,urlcolor=blue,citecolor=black}{hyperref}

\documentclass[pdflatex,sn-basic]{sn-jnl}% Basic style with author-year citations

% Math packages first
\usepackage{amsmath}%
\usepackage{amssymb}%
\usepackage{amsfonts}%
\usepackage{amsthm}%
\usepackage{amsopn}%
\usepackage{amstext}%
\usepackage{mathrsfs}%
\usepackage{stmaryrd}% For llbracket, rrbracket
\usepackage{mathpartir}%

% Graphics and tables
\usepackage{graphicx}%
\DeclareGraphicsExtensions{.png,.pdf,.jpg}% PNG first, then PDF
\usepackage{subcaption}% For subfigure environment
\usepackage{caption}% For caption customization
\captionsetup{justification=raggedright, singlelinecheck=false}% Left-align all captions by default
\captionsetup[figure]{justification=centering, singlelinecheck=false}% Center-align figure captions
\usepackage{tikz}%
\usepackage{multirow}%
\usepackage{booktabs}%
\usepackage{tabularx}
\usepackage{array}

% Algorithm packages
\usepackage{algorithm}%
\usepackage{algorithmicx}%
\usepackage{algpseudocode}%

% Other packages
\usepackage{float}% For [H] option to force exact placement
\usepackage{xcolor}%
\usepackage{listings}%
\usepackage{comment}
\usepackage[authoryear,round]{natbib}% Author-year citation style
\newcolumntype{L}{>{\raggedright\arraybackslash}X}

% Define Solidity language for listings
\lstdefinelanguage{Solidity}{
  keywords={contract, function, modifier, public, private, external, internal, view, pure, payable, returns, return, if, else, for, while, require, assert, mapping, address, uint256, uint, bool, string, memory, storage, calldata, event, emit, struct, enum, constructor, msg, this, new, delete, break, continue, throw, revert},
  keywordstyle=\color{blue}\bfseries,
  ndkeywords={true, false, null, this, new},
  ndkeywordstyle=\color{darkgray}\bfseries,
  identifierstyle=\color{black},
  sensitive=true,
  comment=[l]{//},
  morecomment=[s]{/*}{*/},
  commentstyle=\color{purple}\ttfamily,
  stringstyle=\color{red}\ttfamily,
  morestring=[b]',
  morestring=[b]"
}

% Listings default settings
\lstset{
  language=Solidity,
  basicstyle=\footnotesize\ttfamily,
  backgroundcolor=\color[RGB]{248,248,248},
  numbers=left,
  numbersep=6pt
}

% Define postbreak symbol for listings
\newcommand{\lstcontinuemark}{\mbox{\textcolor{gray}{\(\hookrightarrow\)}}\space}

\theoremstyle{thmstyleone}%
\newtheorem{theorem}{Theorem}%  meant for continuous numbers
\newtheorem{proposition}[theorem]{Proposition}%
\theoremstyle{thmstyletwo}%
\newtheorem{example}{Example}%
\newtheorem{remark}{Remark}%

\theoremstyle{thmstylethree}%
\newtheorem{definition}{Definition}%

\raggedbottom

% Using author-year style to match submission guidelines
\bibpunct{(}{)}{;}{a}{,}{ } % a = author-year citations

% URL breaking settings (hyperref and url already loaded by sn-jnl.cls)
\Urlmuskip=0mu plus 1mu
\def\UrlBreaks{\do\/\do-\do\_\do\.\do\?\do\&}

\makeatletter
\renewcommand\@biblabel[1]{} % Remove numeric labels like [1]
\makeatother
\setlength{\bibsep}{0.2em}   % Space between bibliography items
\setlength{\bibhang}{2em}

% Custom math commands
\newcommand{\Abort}{\mathsf{Abort}}
\newcommand{\Norm}{\mathsf{Norm}}
\newcommand{\Ret}{\mathsf{Ret}}
\newcommand{\Env}{\mathit{Env}}


\begin{document}

\title[Solidity Guardian]{Solidity Guardian: Validating Developer Intentions in Smart Contracts via Intent Annotations to Mitigate Numeric Logical Vulnerability}

\author[1]{\fnm{Inseong} \sur{Jeon}}\email{iwwyou@korea.ac.kr}

\author[1]{\fnm{Sundeuk} \sur{Kim}}\email{sd\_kim@korea.ac.kr}

\author[1]{\fnm{Hyunwoo} \sur{Kim}}\email{khw0809@korea.ac.kr}

\author*[1]{\fnm{Hoh Peter} \sur{In}}\email{hoh\_in@korea.ac.kr}

\affil*[1]{\orgdiv{Department of Computer Science}, \orgname{Korea University},
\orgaddress{\street{145, Anam-ro}, \city{Seonbuk-gu}, \postcode{02841}, \state{Seoul},
\country{Republic of Korea}}}


\abstract{In smart contract development, arithmetic issues such as integer division truncation, incorrect operation order, and scaling mismatches can silently violate the developer's intended numeric relationships—yet still result in successful transactions. These \textit{numeric logical vulnerabilities} typically evade runtime exceptions, static analysis rules, and dynamic testing, often leading to significant financial losses after deployment. This paper introduces \textbf{Solidity Guardian}, a lightweight tool that enables developers to declare intended numeric invariants directly in code and receive immediate feedback on compliance. By leveraging an \textit{Intent DSL}—structured around \texttt{@Pre}, \texttt{@During}, and \texttt{@Post} phases—and combining it with SolQDebug's interval-based abstract interpretation engine, Solidity Guardian continuously tracks variable ranges from function entry to exit. The system performs incremental parsing and selective control-flow analysis, allowing it to deliver deterministic \texttt{True/False/Unknown} verdicts within milliseconds (average 30 ms for functions up to 500 lines), seamlessly integrating with the development workflow. Our results demonstrate that this "intent annotation + deterministic checking" approach effectively prevents subtle numeric logic errors during development, proposing a new paradigm for mechanically verifying developer intent in smart contract programming.}


\keywords{Smart Contract Development, Solidity, Numeric Vulnerabilities, Intent Annotations, Abstract Interpretation, Developer Tools}

\maketitle
\section{Introduction}
% 여기서 왜 하필 Numeric에 집중했는지 써야됨 - AI 때문에, 값분석이 중요한 부분에서는 AI보다는... 그리고 디버깅 관점에서도

Smart contracts handle assets directly by performing arithmetic operations on integer values such as token balances, fees, and percentages.  
While Solidity 0.8 and later versions enforce automatic overflow checks, logical numeric errors—such as incorrect operation order, integer division truncation, and scaling mismatches—remain elusive and are not easily detected by conventional means.

Although large language models (LLMs) have recently been introduced to support tasks such as code completion, refactoring, and bug explanation, they often struggle with precise value inference and boundary condition analysis due to their generative nature.  
Crucially, any financial loss resulting from LLM-generated code ultimately becomes the developer’s responsibility, underscoring the need for rigorous validation of numeric invariants during development.  
A notable real-world example is the 2023 Sperax DAO incident.  
A single-line bug in the \texttt{\_balanceOf} function caused account balances to be massively overstated.  
Since the transaction executed successfully, monitoring systems failed to detect the issue until after approximately \$300,000 worth of tokens had been exploited.

Existing tools often miss such vulnerabilities.  
Static analyzers excel at detecting predefined patterns (e.g., unsafe \texttt{tx.origin} use or unchecked external calls) but are not designed to infer specific numeric intentions like “distribute 10\% of reward.”  
Fuzzers and symbolic execution engines may fail to hit critical boundary conditions, especially when confronted with path explosion in complex mappings, loops, or modifiers.  
Comment-code inconsistency (CCI) detectors provide probabilistic assessments based on weak correlations between natural language comments and code but cannot deliver definitive True/False verification.  
Formal verification offers precision but demands exhaustive specification efforts, rendering it impractical for most real-world projects.  
Thus, there is a pressing need for a lightweight mechanism that allows developers to express numeric intent alongside code and instantly verify its correctness.

To fill this gap, we propose Solidity Guardian, a tool that enables developers to declare intent with a single-line annotation adjacent to the code.  
Our domain-specific language (DSL) adopts a three-phase model—\texttt{@Pre}, \texttt{@During}, and \texttt{@Post}—to specify initial conditions, runtime constraints, and postconditions, respectively. Solidity Guardian leverages SolQDebug's interval-based abstract interpretation engine to compute variable ranges in the background and checks them against the declared intent.  
The system performs partial re-analysis only on the changed code fragments, returning compliance results—\texttt{Satisfied}, \texttt{Violated}, or \texttt{Unknown}—within an average of 30 milliseconds for functions up to 500 lines.  
As a result, developers receive real-time feedback via visual cues directly within the IDE, without disrupting their workflow.

This paper makes the following contributions:
\begin{itemize}
    \item \textbf{Intent DSL and Formal Semantics}: We define a lightweight annotation language based on three execution phases and formalize its semantics in terms of interval state transitions.
    \item \textbf{Millisecond-Scale Verification Engine}: By reusing SolQDebug’s dynamic control-flow graph (CFG) and interval analysis, we achieve millisecond-level feedback through partial recomputation.
    \item \textbf{Large-Scale Empirical Evaluation}: Our evaluation on 1,023 DappSCAN functions and 120 real-world bugs reported by CHEN~\textit{et al.} demonstrates a 92.5\% additional detection rate with only a 3.1\% false positive rate.
    \item \textbf{Practical Use Cases}: We analyze scenarios involving tokenomics (e.g., fee distributions, reward sharing, burn limits), staking contracts, and AMM exchange rates to illustrate how Guardian provides immediate and actionable feedback.
\end{itemize}

The rest of this paper is organized as follows.  
We first review Solidity's numeric types and real-world examples of numeric logic flaws, followed by an analysis of why existing tools fail to detect them.  
We then introduce the Intent DSL and its formal semantics, followed by the design of Guardian's verification engine based on partial re-analysis.  
Subsequently, we present our evaluation results on responsiveness and detection effectiveness, discuss practical applications, limitations, and future directions, and conclude with a comparison to related work.


\section{Background}

\begin{table}[!t]
        \centering
        \caption{Distribution of Variable Types in Smart Contract Functions}
        \label{tab:variable_distribution}
        \resizebox{\columnwidth}{!}{
        \begin{tabular}{lcc}
            \hline
            \textbf{Category} & \textbf{Number of Functions} & \textbf{Percentage (\%)} \\
            \hline
            \textbf{Total Functions Analyzed} & 1524 & \\
            Excluded Functions & 501 & \\           
            \textbf{Functions Considered for Analysis} & 1023 & \\
            \hline
            Functions using \textbf{uint} & 798 & 78.00587 \\
            Functions using \textbf{int} & 5 & 0.48876 \\
            Functions using \textbf{uint} and \textbf{int} & 41 & 4.10557 \\             
            Functions using other types & 178 & 17.39980 \\
            \hline
        \end{tabular}}
    \end{table}

\subsection{Solidity and Its Numeric Types}

Solidity is a statically typed, contract-oriented programming language designed for the Ethereum Virtual Machine (EVM)~\citep{solidity, soldom}. Its syntax borrows elements from C++, JavaScript, and Python, and it supports modular development through contracts composed of variables and functions~\citep{soldom, solinf}. As a Turing-complete language, Solidity enables the implementation of complex decentralized applications with persistent on-chain state.

Solidity supports a range of built-in data types~\citep{soltype}:
\begin{itemize}
    \item \textbf{Booleans} — logical \texttt{true} or \texttt{false} values.
    \item \textbf{Integers} — signed (\texttt{intN}) and unsigned (\texttt{uintN}) integer types with widths from 8 to 256 bits.
    \item \textbf{Fixed-point numbers} — declared but currently unsupported in arithmetic operations.
    \item \textbf{Address} — 20-byte externally owned accounts (EOAs) or contract addresses.
    \item \textbf{Fixed-size byte arrays} (\texttt{bytesN}) — statically sized byte sequences.
    \item \textbf{Strings} and \textbf{dynamic bytes} — variable-length UTF-8 strings or byte sequences.
    \item \textbf{Enums} — user-defined enumerated types.
    \item \textbf{User-defined value types} — new type aliases with stricter invariants based on elementary types.
\end{itemize}

Variables of these types can be declared in three storage categories~\citep{solgvar, solsvar}:
\begin{itemize}
    \item \textbf{Global variables} — predefined variables and functions that expose blockchain metadata or serve as utility features (e.g., \texttt{msg}, \texttt{block}, \texttt{tx}).
    \item \textbf{State variables} — persistent contract storage retained across transactions.
    \item \textbf{Local variables} — ephemeral data residing in a function’s execution context.
\end{itemize}

Among these, numeric variables play a dominant role in smart contract logic. To quantify this, we analyzed a DappSCAN-Source dataset~\citep{dappscan} comprising contracts written in Solidity version 0.8.0 or later. From 3,345 available contracts, we randomly selected approximately 10\% of contracts from each size bracket (1–10 KB, 11–20 KB, and over 20 KB), yielding 103 Solidity files containing 128 contracts and 1,524 functions in total. We excluded 501 trivial functions, which consist solely of assignments, simple returns, or void function calls, as they lack substantive logic relevant to intent specification. 

Table.1 summarizes our findings. Among the remaining 1,023 functions, over 82\% use numeric variables, with \texttt{uint} appearing in more than 78\% of cases. The prevalence of unsigned integers reflects the dominant use of numeric logic in token-related operations such as \texttt{mint}, \texttt{transfer}, \texttt{withdraw}, \texttt{burn}, \texttt{stake}, and \texttt{reward}. These functions typically enforce non-negative balances, making \texttt{uint} the preferred type. This empirical evidence highlights that numeric types are pervasive in Solidity contracts, underscoring the importance of carefully handling arithmetic operations during development to prevent critical logic errors.


\subsection{Numeric Logical Vulnerabilities in Smart Contracts}

    \begin{figure}
                \centering
          \begin{lstlisting}[
              language=Solidity,
              basicstyle=\scriptsize\ttfamily,
              columns=flexible,
              breaklines=true,
              breakatwhitespace=false,
              breakindent=0pt,
              postbreak=\lstcontinuemark,
              numbers=left, numbersep=6pt, xleftmargin=1.5em
          ]
// ..Code for other functions are omitted
function _ensureRebasingMigration(address _account) internal {
    if (nonRebasingCreditsPerToken[_account] == 0) {
        nonRebasingCreditsPerToken[_account] = 1;
        if (_creditBalances[_account] != 0) {            
            uint256 bal = _balanceOf(_account);
            nonRebasingSupply = nonRebasingSupply.add(bal);
            _creditBalances[_account] = bal;
        }            
    }
}
function _balanceOf(address _account) private view returns (uint256) {
    uint256 credits = _creditBalances[_account];
    if (credits > 0) {
        if (nonRebasingCreditsPerToken[_account] > 0) {
            // The code should have been 
            // credits/creditPerToken
            return credits;
        }
        return credits.divPrecisely(rebasingCreditsPerToken);
    }
    return 0;
}

// ..Code for other functions are omitted
        \end{lstlisting}
        \caption{Motivation Example~\citep{fig1}}
    \end{figure}
    
In blockchain systems, even transactions that complete successfully at runtime may harbor severe \textit{numeric logical vulnerabilities} if they violate the developer’s intended logic due to subtle errors such as incorrect operation order, integer division truncation, or faulty branching conditions. Unlike runtime exceptions, these logic errors do not trigger reverts or explicit failures. Consequently, they can bypass both conventional testing and on-chain monitoring, becoming permanently recorded on the blockchain.

A notable real-world example is the \textit{Sperax DAO incident} in February 2023. The root cause was a one-line logic error in the \texttt{\_balanceOf} function (Figure~\ref{fig:sperax}), which mistakenly returned the raw \texttt{credits} variable instead of performing the intended normalization \texttt{credits / creditPerToken}. This seemingly minor mistake led to an inflated account balance that exceeded the correct value by several hundred times. The overestimated balance was subsequently aggregated into the total supply via the \texttt{\_ensureRebasingMigration} function. Because all transactions executed successfully without exceptions, on-chain monitors failed to detect the anomaly. By the time the exploit was noticed, attackers had already drained tokens worth approximately \$300,000 under the guise of legitimate withdrawals.

Furthermore, CHEN \textit{et al.}~\citep{numscout} conducted a systematic analysis of 1,199 audit reports from the DappSCAN dataset~\citep{dappscan}, identifying five distinct categories of novel numeric logical vulnerabilities:

\begin{itemize}
    \item \textbf{Div-In-Path}: Integer division within conditional expressions distorts control flow.  
    A condition such as \texttt{if (msg.value / 1 ether > 3)} wrongly evaluates to false for \texttt{3.9 ETH}, failing the branch. Eleven real-world attack cases were documented
    
    \item \textbf{Operator-Order Error}: Arithmetic operations applied in the wrong order, causing premature integer truncation.  
    For example, \texttt{reward / 100 * 10} always evaluates to zero when \texttt{reward < 100}.    
    .
    
    \item \textbf{Minor Retention}: Rounding remainders (e.g., fees or interest) become permanently locked in the contract.  
    In one case, successive remainder accumulation in a distribution loop resulted in 90 ETH remaining locked over three months.
    
    \item \textbf{Exchange Problem}: Scaling factor mismatches cause exchange rate deviations by up to \(10^{-18}\).  
    Four DeFi AMM contracts were observed to exhibit slippage losses exceeding 0.3\%.
    
    \item \textbf{Precision-Loss Trend}: Repeated multiplication and division operations amplify rounding errors over time, potentially driving balances negative in contract end states.  
    This vulnerability forced emergency upgrades in two staking pool contracts.
\end{itemize}

Since Solidity 0.8 introduced automatic integer overflow checks with revert semantics, traditional \textit{integer overflow} vulnerabilities have been largely mitigated. However, the aforementioned logical numeric errors remain undetectable by runtime safeguards, as they manifest only in semantically incorrect—but syntactically valid—execution outcomes. Therefore, unless developers explicitly verify numeric invariants on critical quantities such as amounts, ratios, or balances during development, these vulnerabilities are virtually impossible to detect using conventional methods.


%────────────────────────────────────────────
\subsection{Limitations of Existing Detection Approaches}



Smart contract analysis tools have advanced across various categories, including static analysis, dynamic analysis, comment-code inconsistency (CCI) detection, and formal verification. However, most existing tools struggle to deterministically detect numeric logic errors—cases where the execution result violates the developer’s intent despite syntactically correct code. Consider the following simple yet critical example:

\begin{lstlisting}[language=Solidity]
// Intended: 10% of reward
// Actual: 0 wei (when reward < 100)
devReward = reward.div(100).mul(10);
\end{lstlisting}

If \texttt{reward} is between 1 and 99, the integer division \texttt{reward / 100} rounds down to zero, and the multiplication yields zero regardless of the intended 10\% computation. The compiler emits no warnings, nor does the overflow/underflow checker flag this issue. Why does this happen?

\paragraph{(1) Static Analysis (Pattern- and Rule-Based)}
Tools like \textsc{Slither} and \textsc{Mythril} focus on detecting fixed-form defects such as missing \textsc{SafeMath} checks, common reentrancy patterns, or misuse of \texttt{tx.origin}. However, errors related to arithmetic order, rounding behavior, or control-flow distortion—which require comparing intended versus actual numeric values—are difficult to encode as static patterns or rules. Consequently, expressions like the example above typically pass static analysis as “clean.”

\paragraph{(2) Dynamic Analysis, Fuzzing, and Concolic Execution}
Tools like \textsc{Echidna} and \textsc{SFuzz} explore execution paths through randomized or symbolic inputs. Nevertheless, discovering a critical boundary input (e.g., \texttt{reward = 37}) that triggers a logic violation is inherently difficult. Moreover, functions with complex mappings, loops, or modifiers exacerbate state-space explosion, causing tools to hit depth limits and miss violations.

\paragraph{(3) Comment-Code Inconsistency (CCI) Detection with Natural Language Processing}
Recent LLM- or embedding-based approaches attempt to align natural language comments (e.g., “distribute 10\% fee”) with code semantics. However, these models struggle to \textit{prove} a precise numeric relationship—such as “exactly 10\%”—and instead return only probabilistic mismatch scores. If comments are absent or ambiguous, this gap widens further.

\paragraph{(4) Formal Verification (e.g., KEVM, VeriSmart)}
Formal verification frameworks require developers to write explicit invariants or state-machine specifications. In practice, the steep learning curve and high annotation cost limit their adoption in real-world projects. Any code segments lacking formal specifications remain unanalyzed, creating blind spots.

\smallskip

As a result, numeric logic bugs like the example above are permanently recorded on-chain as \texttt{success} transactions. Developers often only discover the problem after user losses occur. Two critical gaps explain why existing tools fail to catch these issues:
\begin{itemize}
    \item \textbf{Lack of Intent Annotations}: Source code alone provides no way to infer that, for example, a value should represent exactly 10\% of another.
    \item \textbf{Lack of Deterministic Checking}: Even when intent is described—through comments or annotations—existing tools cannot deterministically decide the correctness of a single arithmetic expression at the line level.
\end{itemize}

To fill this gap, we propose \textsc{Solidity Guardian}, a framework that allows developers to declare their intended numeric logic using a single-line domain-specific language (DSL). The system then applies interval-based abstract interpretation to compute and compare actual outcomes within milliseconds, providing immediate \textbf{True/False} feedback. This approach offers deterministic verification at sub-wei precision, addressing a critical need that previous tools have left unfilled.


%────────────────────────────────────────────
\subsection{SolQDebug: Foundation for Guardian}

SolQDebug is a source-level Solidity debugger designed to enable interactive abstract interpretation during code authoring. Rather than relying on traditional deployment and transaction-based testing workflows, SolQDebug performs interval-based static analysis on incomplete code fragments and reports symbolic value ranges directly in the Solidity editor. This millisecond-scale feedback loop allows developers to inspect how symbolic inputs propagate through variables and control structures, without compiling or deploying contracts.

At the core of SolQDebug is an extended Solidity grammar that supports inline annotations for symbolic inputs. Developers declare ranges for global, state, and local variables using one-line directives (e.g., \texttt{// @StateVar x = [0, 100]}). The system incrementally constructs a dynamic control-flow graph (CFG) from live edits and executes abstract interpretation over a fixed-point lattice of interval values. Each edit updates only the affected region, enabling fast recomputation. A test-case isolation mechanism ensures that symbolic inputs do not interfere across runs.

In the context of this work, we extend SolQDebug's analysis core as the verification backend for \textit{Solidity Guardian}. Specifically, the interval domain interpreter, dynamic CFG builder, and per-statement memory model are adopted as-is. However, whereas SolQDebug was designed to compute "what values a variable can take," Guardian checks "whether the observed behavior conforms to the developer's intent."

To support this shift in goal, we augment the interpreter to compare computed intervals against developer-declared postconditions and invariants. These logical intent specifications, written in a new annotation language (Section~4.1), allow Guardian to verify correctness properties such as value preservation, expected directional change, or relational constraints across entry, assignment, and exit. The reuse of SolQDebug's infrastructure ensures that verification is lightweight, precise, and responsive—consistent with Guardian's goal of bridging the gap between program logic and developer intent.


%────────────────────────────────────────────
\subsection{Proposed Work: Solidity Guardian}


To address the two critical gaps identified in the previous section—(1) the inability to infer a developer's numeric intent from code alone and (2) the lack of a mechanism to deterministically check whether a single-line arithmetic expression satisfies that intent—we propose \textbf{Solidity Guardian}. The core idea is straightforward:  
Developers declare numeric invariants, such as ``this value must be 10\%'' or ``\texttt{balance} should be non-negative after the loop,'' using a single-line \textit{Intent DSL} placed directly alongside the code. Each time the code is edited, Solidity Guardian computes the actual range of the relevant variables within milliseconds and compares it against the declared intent. The outcome is deterministically categorized into one of three states—\texttt{True}, \texttt{False}, or \texttt{Unknown}—and immediately displayed in the editor using intuitive visual cues.

\paragraph{Intent DSL Summary}
Solidity Guardian introduces a simple domain-specific language with tokens reflecting three analysis stages:
\begin{itemize}
    \item \textbf{@Pre}: Declares the abstract input ranges (e.g., mapping keys, state variables).
    \item \textbf{@During}: Specifies intermediate state constraints at specific program points (\textit{Assign}, \textit{Before/After}).
    \item \textbf{@Post}: Declares the expected relationships at function exit (\textit{Entry/Exit}, \textit{returnValues}).
\end{itemize}

The DSL supports comparison operators $<, >, \leq, \geq, ==, !=$ with right-hand sides defined as constants, variables, intervals, or linear arithmetic expressions.

\begin{lstlisting}[language=Solidity]
// @Pre    _amount = [1, 99]
// @During devReward (Assign == Current * 10 / 100)
// @Post   returnValues == reward * 10 / 100
\end{lstlisting}

\paragraph{Deterministic Checking Procedure}
Upon saving the file or confirming a line of code, Solidity Guardian re-parses only the modified code fragment. It leverages SolQDebug's verified dynamic control flow graph (CFG) and interval propagation engine to compute the actual value intervals of variables.

\[
  \textbf{Satisfied (✓)} \quad \Longleftrightarrow \quad
    I_{\text{actual}} \subseteq I_{\text{intent}}
\]

\[
  \textbf{Violated (✗)} \quad \Longleftrightarrow \quad
    I_{\text{actual}} \cap I_{\text{intent}} = \varnothing
\]

\[
  \textbf{Uncertain (△)} \quad \Longleftrightarrow \quad
    I_{\text{actual}} \not\subseteq I_{\text{intent}} 
    \;\;\wedge\;\; I_{\text{actual}} \cap I_{\text{intent}} \neq \varnothing
\]

If neither satisfied nor violated, the system conservatively reports \texttt{Unknown}.  
The entire verification process is implemented in a Python backend, consistently completing within an average of 30 milliseconds even for functions around 500 lines long.

\paragraph{Key Advantages}
\begin{itemize}
    \item \textbf{Developer-Centric Intent Declarations}:  
    Numeric invariants can be expressed with a single line of comment-like syntax, reducing the learning curve.
    
    \item \textbf{Millisecond-Level Responsiveness}:  
    Partial recalculations allow feedback with minimal latency, comparable to keystroke delays.
    
    \item \textbf{Deterministic Results}:  
    By providing clear \texttt{True}/\texttt{False} outcomes, the tool reliably detects even sub-wei deviations.
\end{itemize}

In the following sections, we present the formal \textit{Intent Model} used by Solidity Guardian and detail the verification procedure that determines compliance between declared intent and actual execution outcomes.


%════════════════════════════════════════════════════════════════════
\section{The Design of Solidity Guardian}
%════════════════════════════════════════════════════════════════════

%────────────────────────────────────────────────────────
\subsection{System Architecture}
%────────────────────────────────────────────────────────

% TODO: Add system architecture description
% - Overview of Guardian's components
% - Integration with SolQDebug
% - Workflow: annotation → parsing → verification → feedback


%────────────────────────────────────────────────────────
\subsection{Intent Model}
\label{sec:intent-model}
%────────────────────────────────────────────────────────

\paragraph{Goal.}
\textsc{Solidity Guardian} provides a light-weight, single-line
annotation language that lets developers state—next to the code—
\emph{what should hold}, not how to prove it.
Each annotation is checked incrementally after every edit against
an interval abstract state, producing a deterministic
\texttt{True/False/Unknown} in a few milliseconds, plus a numeric
\emph{confidence} and, when useful, a \emph{diagnostic interval}
explaining where the uncertainty comes from.

\vspace{0.25em}
\noindent\textbf{Top-level shape.}
An annotated file is a bag of \emph{intent directives} partitioned
into three phases:
\[
\texttt{@Pre}\quad\texttt{@During}\quad\texttt{@Post}
\]
They may appear in any textual order but are interpreted in the
semantic order:
\emph{seed injection} $\rightarrow$ \emph{mid-execution checks}
$\rightarrow$ \emph{exit-time checks}.
The entry non-terminal is \texttt{intentUnit}.

\begin{lstlisting}[basicstyle=\scriptsize\ttfamily,
                   numbers=left, frame=single,
                   caption={Guardian annotation grammar (top-level excerpt)},
                   label={lst:intent-grammar-top}]
intentUnit
  : ( preDirective | duringDirective | postDirective )* EOF ;

preDirective
  : '//' '@GlobalVar' identifier ('.' identifier)? '=' seedValue  # PreGlobal
  | '//' '@StateVar'  varRef '=' seedValue                        # PreState
  | '//' '@LocalVar'  varRef '=' seedValue                        # PreLocal ;

duringDirective
  : '//' '@During' duringFormula (',' duringFormula)* ;
postDirective
  : '//' '@Post'   postFormula   (',' postFormula  )* ;
\end{lstlisting}

\paragraph{Formulas and clauses.}
A \emph{formula} is a disjunction/conjunction of \emph{clauses}.
Each clause is either a primitive predicate or an implication
(“if–then”) between two predicates.  Logical operators use the
usual precedence and are left-associative; short-circuiting is
observed in the checker.

\begin{lstlisting}[basicstyle=\scriptsize\ttfamily,frame=single]
duringFormula : duringClause (logicOp duringClause)* ;
postFormula   : postClause   (logicOp postClause  )* ;
duringClause  : predicateDuring
              | predicateDuring '=>' predicateDuring  # DuringImplication ;
postClause    : predicatePost
              | predicatePost   '=>' predicatePost    # PostImplication ;
logicOp       : '&&' | '||' ;
\end{lstlisting}

%────────────────────────────────────────────
\paragraph{Primitive predicates.}
We distinguish \emph{phase-specific} predicates (temporal snapshots)
from \emph{common} predicates that are shared across phases.

\begin{lstlisting}[basicstyle=\scriptsize\ttfamily,frame=single]
predicateDuring
  : varRef '(' 'Before'  relOp 'After'   ')'  # DuringBeforeAfter
  | varRef '(' 'Assign'  relOp 'Current' ')'  # DuringAssignCurrent
  | commonPredicate                           # DuringCommonPredicate ;

predicatePost
  : varRef '(' 'Entry' relOp 'Exit' ')'       # PostEntryExit
  | 'Unchanged' '(' varRef ')'                # UnchangedVar
  | commonPredicate                           # PostCommonPredicate ;

commonPredicate
  : 'returnExpression' relOp value            # ReturnExprCmp
  | 'return' varRef      relOp value          # ReturnVarCmp
  | value                relOp value          # RelationalCmp ;

relOp : '<' | '>' | '<=' | '>=' | '==' | '!=' | 'in' | 'not' 'in' ;
\end{lstlisting}

\noindent
\textbf{Snapshot predicates (single-variable only).}
To avoid ambiguity and encourage local reasoning, the snapshot
forms intentionally restrict the left-hand side to a single
\texttt{varRef}.  The operator is written \emph{between} named
snapshots:
\begin{center}
\verb|v(Before < After)|,\quad
\verb|v(Assign == Current)|,\quad
\verb|v(Entry <= Exit)|.
\end{center}
Here, \texttt{Assign} captures the value immediately after the
current assignment to \texttt{v}; \texttt{Before}/\texttt{After}
denote the state around the current statement; and
\texttt{Entry}/\texttt{Exit} refer to function entry/exit.

\noindent
\textbf{Return predicates.}
\verb|returnExpression| refers to the unnamed return value at the
relevant program point (\texttt{@During}: the particular \texttt{return}
statement; \texttt{@Post}: the join over normal exits).
\verb|return v| projects the $i$-th element for tuple returns via
\verb|return[i]|.

\noindent
\textbf{Free comparisons.}
Any two \emph{intent expressions} (\S\ref{sec:intent-expr}) can be
compared using a relational operator, including membership
\verb|in| / \verb|not in|.
We normalise \verb|not in| lexically (whitespace-insensitive).

%────────────────────────────────────────────
\paragraph{Implication (\texorpdfstring{$\Rightarrow$}).}
A clause \verb|A => B| reads “if A holds, then B must hold”.
We use three-valued, short-circuit semantics
(\textsf{T}, \textsf{F}, \textsf{U} for Unknown):
\[
\begin{array}{c|ccc}
A\Rightarrow B & \textsf{T} & \textsf{F} & \textsf{U} \\\hline
\textsf{T}      & \textsf{T} & \textsf{F} & \textsf{U}\\
\textsf{F}      & \multicolumn{3}{c}{\textsf{T} \text{ (vacuously true)}} \\
\textsf{U}      & \multicolumn{3}{c}{\textsf{U}}
\end{array}
\]
That is, a false antecedent makes the implication hold
vacuously; an unknown antecedent yields an unknown implication.

%────────────────────────────────────────────
\subsubsection{Intent expressions}
\label{sec:intent-expr}
%
Every operand in \texttt{@During} and \texttt{@Post} is an
\emph{intent expression} drawn from three families; the grammar
below also serves as the developer-facing syntax.

\begin{lstlisting}[basicstyle=\scriptsize\ttfamily,frame=single]
value      : arithExpr | addressExpr | boolExpr ;
arithExpr  : arithExpr ('+'|'-') arithTerm | arithTerm ;
arithTerm  : arithTerm ('*'|'/'|'%') arithFactor | arithFactor ;
arithFactor: signedNumberLiteral
           | '[' signedNumberLiteral ',' signedNumberLiteral ']'
           | varRef
           | 'PercentOf' '(' arithExpr ',' numberLiteral ')'
           | 'ceil'  '(' arithExpr ',' numberLiteral ')'
           | 'floor' '(' arithExpr ',' numberLiteral ')'
           | '(' arithExpr ')' ;
addressExpr: 'address' numberLiteral
           | 'symbolicAddress' numberLiteral
           | varRef ;
boolExpr   : booleanLiteral | varRef ;
\end{lstlisting}

\noindent
\textbf{Helpers.}
\texttt{PercentOf$(x,p)$}, \texttt{ceil$(n,d)$}, \texttt{floor$(n,d)$}
are desugared arithmetically prior to evaluation.  This keeps the
checker simple: it only needs interval arithmetic.

%────────────────────────────────────────────
\subsubsection{Pre‑Intent: seeding the abstract store}
%
\texttt{@Pre} directives inject symbolic seeds before analysis:
intervals, Booleans, fixed/symbolic addresses, enums, and inline
arrays.  They can target globals, state variables, or locals.


%────────────────────────────────────────────────────────
\subsection{Guardian Verification Engine}
%────────────────────────────────────────────────────────

% TODO: Add verification engine description
% - How interval-based checking works
% - Comparison between actual intervals and intent specifications
% - Three-valued logic (Satisfied, Violated, Unknown)
% - Performance optimizations (partial re-analysis)


%────────────────────────────────────────────────────────
\subsection{Running Example}
%────────────────────────────────────────────────────────

% TODO: Add running example
% - Show one or more of NumScout's 5 defect types
% - Demonstrate how to specify intent with Guardian DSL
% - Show verification results
% Example candidates:
% - Operator-Order Error: reward / 100 * 10 → 0 when reward < 100
% - Div-In-Path: if (msg.value / 1 ether > 3) fails for 3.9 ETH
% - Minor Retention: fee remainders locked in contract


%════════════════════════════════════════════════════════════════════
\section{Evaluation}
%════════════════════════════════════════════════════════════════════

\subsection{Experimental Setup}

% TODO: Add experimental setup description
% - Dataset: NumScout 95 labeled samples + DappSCAN functions
% - Baseline: NumScout
% - Metrics: Detection rate, false positives, response time
% - Implementation details


\subsection{RQ1: Effectiveness in Detecting Known Numeric Defects}

% TODO: Add RQ1 evaluation
% - NumScout의 5가지 defect types를 DSL로 표현
% - 95개 샘플에 대한 detection 결과 비교
% - Guardian으로 표현 가능한 패턴 범위


\subsection{RQ2: Expressiveness of the Guardian DSL}

% TODO: Add RQ2 evaluation
% - Guardian DSL로 NumScout 패턮 외 추가 속성 표현 가능?
% - DSL specification 작성 난이도, 라인 수 등
% - Case studies with different defect patterns


\subsection{RQ3: Performance and Response Time}

% TODO: Add RQ3 evaluation
% - 분석 시간, 메모리 사용량
% - Contract 크기/복잡도에 따른 성능
% - Comparison with NumScout and other tools


%════════════════════════════════════════════════════════════════════
\section{Discussion}
%════════════════════════════════════════════════════════════════════

\subsection{Generalizability: Beyond Numeric Defects}

% TODO: Add discussion on generalizability
% - NumScout의 패턴은 Guardian의 특수 케이스
% - DSL의 확장 가능성
% - Non-numeric logical properties (access control, state consistency 등)
% - "Numeric에 집중한 이유는 명확한 검증 가능성과 높은 impact"


\subsection{Limitations}

% TODO: Add limitations
% - Multi-contract analysis (currently library only)
% - Multi-transaction scenarios (single transaction only)
% - Scalability concerns for very large contracts
% - Need for developer effort in writing intent specifications


%════════════════════════════════════════════════════════════════════
\section{Related Works}
%════════════════════════════════════════════════════════════════════

% 방법론 중심 구조: "개발자 의도를 어떻게 명시하고 검증하는가"
% 세 가지 접근법:
% 1. Pattern-based: 전문가가 정의한 패턴 매칭 (NumScout, Slither)
% 2. NL-based: 자연어 주석에서 의도 추론 (DonCon, SmartCoCo)
% 3. Formal-based: 형식 명세 언어로 의도 표현 (Dafny, Certora)
%
% Guardian의 위치: 세 방법론의 중간
% - Pattern보다 유연 (user-defined intent)
% - NL보다 정확 (structured DSL)
% - Formal보다 실용적 (lightweight, numeric-focused)

\subsection{Pattern-based Vulnerability Detection}

% TODO: Add related works on pattern-based approaches
%
% 핵심 특징: 전문가가 미리 정의한 패턴을 코드에서 찾음
% 장점: 자동화, 빠름, 검증된 패턴
% 단점: 패턴에 없으면 탐지 불가, 확장성 제한, 개발자 의도 표현 불가
%
% Numeric-focused:
% - NumScout~\cite{numscout}: 5가지 numeric defect patterns
%   (Div-In-Path, Operator-Order, Minor Retention, Exchange, Precision-Loss)
%   → LLM-pruning symbolic execution
%   → Hard-coded patterns, 새로운 invariant 추가 불가
%   → Guardian의 핵심 baseline
%
% General Smart Contract Analysis:
% - Slither~\cite{slither}: 70+ detectors, pattern matching
%   → reentrancy, tx.origin, uninitialized storage 등
%   → numeric logic invariants 표현 불가
% - Mythril~\cite{mythril}: symbolic execution + security patterns
%   → integer overflow (pre-0.8), unchecked calls
%   → custom numeric invariants 지원 안 함
% - Securify~\cite{securify}: dataflow patterns in Datalog
%   → compliance/violation patterns
%   → developer intent 명시 메커니즘 없음
% - SmartCheck: XML-based pattern matching on Solidity AST
%
% 차별점:
% - Pattern-based: 도구 개발자가 패턴 정의
% - Guardian: 코드 개발자가 intent 정의
% - NumScout 패턴을 DSL로 표현 가능 + 새로운 invariant 추가 가능
% - Deterministic verification with millisecond feedback


\subsection{Natural Language-based Intent Detection}

% TODO: Add related works on NL-based approaches
%
% 핵심 특징: 자연어 주석/문서에서 의도를 추론하고 코드와 비교
% 장점: 기존 주석 활용, 추가 노력 불필요
% 단점: 주석 품질 의존, 확률적, ambiguous, numeric invariants 표현 어려움
%
% Smart Contract 분석:
% - DonCon~\cite{solapi}: Solidity library API documentation 일치성
%   → function calls, inheritance, events, require/revert
%   → NatSpec comments ↔ AST facts (Datalog)
% - SmartCoCo~\cite{coco}: 주석 ↔ access control, parameter scope
%   → modifiers, zero address checks, require statements
%   → constraint propagation and binding
%
% General Software:
% - @tComment~\cite{tcomment}: Javadoc ↔ code behavior (테스팅 기반)
%   → method parameters, return values, exceptions
% - Panthaplackel et al.~\cite{deep}: 딥러닝 기반 불일치 탐지
%   → code-comment embedding similarity
%   → interactive analysis, edit suggestions
%
% 특징:
% - Natural language로 intent 표현 (유연하지만 ambiguous)
% - 확률적/휴리스틱 기반 불일치 탐지
% - Function-level 분석 중심
% - Numeric invariants 표현 어려움 (e.g., "exactly 10%" → 코드 확인 불가)
%
% 차별점:
% - NL-based: ambiguous, probabilistic, function-level
% - Guardian: structured DSL, deterministic, variable/statement-level
% - @Pre/@During/@Post 3-phase model로 execution trace 명시
% - Numeric invariants에 특화된 표현력 (PercentOf, ceil, floor)
% - True/False/Unknown 판정 (not similarity score)


\subsection{Formal Specification-based Verification}

% TODO: Add related works on formal specification approaches
%
% 핵심 특징: 형식 명세 언어로 correctness property 정의, 정리 증명/SMT solver로 검증
% 장점: 완전한 correctness proof 가능, 엄밀함
% 단점: steep learning curve, 전문 지식 필요, 작성 비용 높음, numeric-specific 아님
%
% General-purpose Formal Verification:
% - Dafny: pre/postconditions, loop invariants, function contracts
%   → decreases clauses, ghost variables
%   → 완전한 correctness proof, 하지만 전문 지식 필요
% - Why3: ML-style specification + multiple SMT backends
%   → Coq, Alt-Ergo, Z3 등 연동
% - Frama-C ACSL: C 프로그램 annotation language
%   → function contracts, assertions, loop invariants
%   → value analysis plugin (interval domain)
%   → 강력하지만 복잡한 specification
%
% Smart Contract Formal Verification:
% - Certora: temporal logic 기반 property specification
%   → invariants, rules, parametric specifications
%   → 강력한 표현력, 하지만 복잡한 CVL 언어
%   → 전문가 수준의 formal methods 지식 요구
% - KEVM: K framework 기반 EVM semantics
%   → bytecode-level formal verification
%   → exhaustive, 하지만 접근성 매우 낮음
% - VeriSmart~\cite{verismart}: transaction-level specification
%   → arithmetic overflow, reentrancy 등
%   → bytecode analysis
% - Zeus~\cite{zeus}: policy specification + abstract interpretation
%   → fairness, data flow policies (XACML-style)
%   → general policies, not numeric intent specific
% - Scribble: runtime assertion language
%   → Solidity annotations → instrumented runtime checks
%   → require/assert 형태로 변환
%   → runtime overhead, not static verification
%
% 차별점:
% - Formal-based: 완전한 specification, 엄밀하지만 복잡함
% - Guardian: lightweight intent DSL (single-line annotation)
% - 학습 곡선 낮음, numeric invariants에 특화
% - Millisecond-scale static verification (vs runtime checks/SMT solving)
% - Deterministic feedback during development (not post-hoc verification)
%
% Technical Foundation (Abstract Interpretation):
% Guardian은 SolQDebug의 interval-based abstract interpretation을 활용
% - Cousot & Cousot (1977): abstract interpretation 이론
% - SolQDebug: Solidity용 interactive interval analysis
%   → "변수 range 계산" → Guardian에서 "intent 검증"으로 확장
% - Astrée, Frama-C: embedded C에 대한 interval analysis


%════════════════════════════════════════════════════════════════════
\section{Conclusion}
%════════════════════════════════════════════════════════════════════


% ─── ① 프리앰블에 토글 스위치 정의 ─────────────────────────
\newif\ifshowack          % 새 불리언 생성
\showackfalse             % false = 숨김, true = 표시
% \showacktrue            % 카메라-레디 때는 이 줄만 주석 해제

% ─── ② 본문에서 조건부로 삽입 ──────────────────────────────
\ifshowack
  \section*{Acknowledgment}  
\fi

\begin{thebibliography}{99}

\bibitem[BlockSecTeam(2023)]{fig1} BlockSecTeam: Sperax USDs: How to Discover Hidden Bug and 300k+ Profit. \url{https://blocksecteam.medium.com/}. Accessed 27 September 2025

\bibitem[Chen et al.(2025)]{numscout} Chen, J., Xia, X., Lo, D., Grundy, J., Yang, X.: NumScout: Unveiling Numerical Defects in Smart Contracts using LLM-Pruning Symbolic Execution. IEEE Trans. Softw. Eng. (2025)

\bibitem[Grieco et al.(2020)]{echidna} Grieco, G., Song, W., Cygan, A., Feist, J., Groce, A.: Echidna: effective, usable, and fast fuzzing for smart contracts. In: Proc. 29th ACM SIGSOFT Int. Symp. Softw. Test. Anal., pp. 557--560 (2020)

\bibitem[Hao et al.(2023)]{coco} Hao, S., Qin, Z., Wang, X., Shu, B., Qin, H., Pan, S.: SmartCoCo: Checking Comment-Code Inconsistency in Smart Contracts via Constraint Propagation and Binding. In: 38th IEEE/ACM Int. Conf. Autom. Softw. Eng. (ASE), pp. 294--306. IEEE (2023)

\bibitem[Hu et al.(2023)]{solidet} Hu, T., Wang, S., Huang, J., Cao, Q., Hu, B., Chen, J.: Detect defects of solidity smart contract based on the knowledge graph. IEEE Trans. Reliab. \textbf{73}(1), 186--202 (2023)

\bibitem[Jiang et al.(2018)]{confuzz} Jiang, B., Liu, Y., Chan, W.K.: Contractfuzzer: Fuzzing smart contracts for vulnerability detection. In: Proc. 33rd ACM/IEEE Int. Conf. Autom. Softw. Eng., pp. 259--269 (2018)

\bibitem[Kalra et al.(2018)]{zeus} Kalra, S., Goel, S., Dhawan, M., Sharma, S.: ZEUS: Analyzing safety of smart contracts. In: Proc. 2018 Netw. Distrib. Syst. Secur. Symp., pp. 1--15. Internet Society (2018)

\bibitem[Language Influences(2025)]{solinf} Language Influences. \url{https://docs.soliditylang.org/en/v0.8.30/language-influences.html}. Accessed 27 September 2025

\bibitem[Ma et al.(2023)]{transracer} Ma, C., Song, W., Huang, J.: TransRacer: Function Dependence-Guided Transaction Race Detection for Smart Contracts. In: Proc. 31st ACM Joint Eur. Softw. Eng. Conf. Symp. Found. Softw. Eng., pp. 947--959 (2023)

\bibitem[Nguyen et al.(2020)]{sfuzz} Nguyen, T.D., Pham, L.H., Sun, J., Lin, Y., Minh, Q.T.: sfuzz: An efficient adaptive fuzzer for solidity smart contracts. In: Proc. ACM/IEEE 42nd Int. Conf. Softw. Eng., pp. 778--788 (2020)

\bibitem[Panthaplackel et al.(2021)]{deep} Panthaplackel, S., Nie, J.Y., Gligoric, M., Li, J.J., Mooney, R.J.: Deep just-in-time inconsistency detection between comments and source code. In: Proc. AAAI Conf. Artif. Intell., pp. 427--435 (2021)

\bibitem[Pasqua et al.(2023)]{ethersolve} Pasqua, M., Ferrara, P., Cortesi, A.: Enhancing Ethereum smart-contracts static analysis by computing a precise Control-Flow Graph of Ethereum bytecode. J. Syst. Softw. \textbf{200}, 111653 (2023)

\bibitem[Shou et al.(2023)]{ityfuzz} Shou, C., Tan, S., Sen, K.: Ityfuzz: Snapshot-based fuzzer for smart contract. In: Proc. 32nd ACM SIGSOFT Int. Symp. Softw. Test. Anal., pp. 322--333 (2023)

\bibitem[So et al.(2020)]{verismart} So, S., Lee, M., Park, J., Lee, H., Oh, H.: Verismart: A highly precise safety verifier for ethereum smart contracts. In: 2020 IEEE Symp. Secur. Priv. (SP), pp. 1678--1694. IEEE (2020)

\bibitem[Solidity Documentation(2025)]{solidity} Solidity Documentation. \url{https://docs.soliditylang.org/en/v0.8.30/}. Accessed 27 September 2025

\bibitem[Stephens et al.(2021)]{pulse} Stephens, J., Ferles, K., Marber, B., Sakkas, G., Dillig, I.: SmartPulse: automated checking of temporal properties in smart contracts. In: 2021 IEEE Symp. Secur. Priv. (SP), pp. 555--571. IEEE (2021)

\bibitem[Structure of a Contract(2025)]{solsvar} Structure of a Contract. \url{https://docs.soliditylang.org/en/v0.8.30/structure-of-a-contract.html}. Accessed 27 September 2025

\bibitem[Tan et al.(2012)]{tcomment} Tan, S.H., Marinov, D., Tan, L., Leavens, G.T.: @tcomment: Testing javadoc comments to detect comment-code inconsistencies. In: 2012 IEEE 5th Int. Conf. Softw. Test. Verif. Valid., pp. 260--269. IEEE (2012)

\bibitem[Tsankov et al.(2018)]{securify} Tsankov, P., Dan, A., Drachsler-Cohen, D., Gervais, A., B{\"u}nzli, F., Vechev, M.: Securify: Practical security analysis of smart contracts. In: Proc. 2018 ACM SIGSAC Conf. Comput. Commun. Secur., pp. 67--82 (2018)

\bibitem[Feist et al.(2019)]{slither} Feist, J., Grieco, G., Groce, A.: Slither: a static analysis framework for smart contracts. In: 2019 IEEE/ACM 2nd Int. Workshop Emerg. Trends Softw. Eng. Blockchain (WETSEB), pp. 8--15. IEEE (2019)

\bibitem[Types(2025)]{soltype} Types. \url{https://docs.soliditylang.org/en/v0.8.30/types.html}. Accessed 27 September 2025

\bibitem[Units and Globally Available Variables(2025)]{solgvar} Units and Globally Available Variables. \url{https://docs.soliditylang.org/en/v0.8.30/units-and-global-variables.html}. Accessed 27 September 2025

\bibitem[Vidal et al.(2024)]{detsurvey} Vidal, F.R., Ivaki, N., Laranjeiro, N.: Vulnerability detection techniques for smart contracts: A systematic literature review. J. Syst. Softw. \textbf{212}, 112160 (2024)

\bibitem[Wohrer and Zdun(2018)]{soldom} Wohrer, M., Zdun, U.: Smart contracts: security patterns in the ethereum ecosystem and solidity. In: 2018 Int. Workshop Blockchain Orient. Softw. Eng. (IWBOSE), pp. 2--8. IEEE (2018)

\bibitem[Wu et al.(2021)]{peculiar} Wu, H., Zhang, Z., Wang, S., Lei, Y., Lin, B., Qin, Y., Zhang, H., Gui, X.: Peculiar: Smart contract vulnerability detection based on crucial data flow graph and pre-training techniques. In: 2021 IEEE 32nd Int. Symp. Softw. Reliab. Eng. (ISSRE), pp. 378--389. IEEE (2021)

\bibitem[Yao et al.(2022)]{mythril} Yao, Y., Liu, P., Gong, H., Li, F., Song, X., Guo, F., Chen, X.: An improved vulnerability detection system of smart contracts based on symbolic execution. In: 2022 IEEE Int. Conf. Big Data (Big Data), pp. 3225--3234. IEEE (2022)

\bibitem[Yu et al.(2023)]{pscvfinder} Yu, L., Wang, K., Liu, B., Li, T., Cheng, X., Fang, C., Chen, Z.: Pscvfinder: a prompt-tuning based framework for smart contract vulnerability detection. In: 2023 IEEE 34th Int. Symp. Softw. Reliab. Eng. (ISSRE), pp. 556--567. IEEE (2023)

\bibitem[Zheng et al.(2024)]{dappscan} Zheng, Z., Chen, J., Nan, Y., Zheng, X., Wu, N., Lou, Y., Xing, Z., Xu, X.: Dappscan: building large-scale datasets for smart contract weaknesses in dapp projects. IEEE Trans. Softw. Eng. \textbf{50}(6), 1360--1373 (2024)

\bibitem[Zhou et al.(2022)]{tmlvd} Zhou, Q., Zhong, R., Qin, L., He, L., Yang, K., Tian, C., Li, Z.: Vulnerability analysis of smart contract for blockchain-based IoT applications: a machine learning approach. IEEE Internet Things J. \textbf{9}(24), 24695--24707 (2022)

\bibitem[Zhu et al.(2022)]{solapi} Zhu, C., Miao, Y., Gao, K., Chen, J., Guo, Y., Liu, Y., Quan, W.: Identifying solidity smart contract api documentation errors. In: Proc. 37th IEEE/ACM Int. Conf. Autom. Softw. Eng. (2022)

\bibitem[Zou et al.(2019)]{debugsurvey} Zou, W., Lo, D., Kochhar, P.S., Le, X.B.D., Xia, X., Feng, Y., Chen, Z., Xu, B.: Smart contract development: Challenges and opportunities. IEEE Trans. Softw. Eng. \textbf{47}(10), 2084--2106 (2019)

\end{thebibliography}



\end{document}